% ============================================================
%  abstract.tex — TH-GNN paper
%  Note: \begin{abstract}...\end{abstract} is provided by main.tex
%        (required before \maketitle in acmart).
% ============================================================

Shilling attacks on recommender systems have evolved from simple heuristic injections to sophisticated LLM-agent-generated profiles that mimic genuine user behaviour at scale. Existing detectors fall into two families — feature-based statistical methods and graph neural network approaches — yet neither has been systematically evaluated against the full spectrum of modern attack sophistication. We present TH-GNN, a Temporal Heterogeneous Graph Neural Network that combines heterogeneous graph attention, relative temporal encoding, and multimodal review-text fusion, and deploy it as an empirical probe across four attack families of increasing sophistication: heuristic (Random, Bandwagon), GAN-based (AUSH), and LLM-agent (Agent4SR). Our study reveals a critical asymmetry: feature-based detectors achieve F1 > 0.93 against heuristic attacks but collapse to F1 = 0.31 against GAN-based AUSH profiles, while TH-GNN maintains consistent recall (DR > 0.99) across all attack families at the cost of elevated false alarm rates. Counterintuitively, feature-based detectors remain effective against the current Agent4SR implementation (best F1\,=\,0.86), primarily because LLM-generated profiles still exhibit detectable volume and temporal-concentration artifacts under the experimental setup; graph-based TH-GNN suffers precision collapse under the resulting distribution shift (F1\,=\,0.035). The open frontier is detection when these incidental artifacts are removed and when graph-based detectors can be calibrated without label access.
Code is available at \url{https://github.com/shiv3589/thgnn-shilling}.
